\section{Spectral silence and geometric correlations in Laughlin manifolds}
\label{sec:laughlin-results}

\subsection{A degenerate spectrum with a resolved projector response}

We begin with the obstruction that motivates the geometric calculation.  In
the registered spectral-silence instance, the zero-mode projector has rank
$D=\SpectralFiberRank$.  Its numerical bandwidth is
$\SpectralKernelBandwidth$, whereas the first state outside the fiber lies at
the positive gap $\SpectralExternalGap$.  The former is a residual audit and
the latter is the scale that permits a reduced resolvent; neither is obtained
by unfolding a narrow band.  Within the accepted fiber all energies are
exactly equal at the algebraic parent, up to the reported solver residual.

The energy form factor consequently contains no information about the
orientation of $P$.  With the normalization used in Fig.~\ref{fig:spectral},
the raw contribution is the constant rank $D$, and the connected energy
spectral form factor vanishes.  This remains true for every tangent panel
because the energy calculation uses only the repeated eigenvalue.  It does not
imply a vanishing response: the same projector has nonzero
$X_a=Q(\partial_aP)P$ under the registered constraint or local-potential
tangents.

Figure~\ref{fig:spectral}(b) makes this separation visible without invoking a
single scalar summary.  A structured curvature, a physical curvature panel,
and a finite-$D$ Jacobi reference can have the same matrix rank while having
clearly different eigenvalue configurations.  The physical spectrum fills
the allowed interval and shows local repulsion, whereas the structured
control occupies a much smaller set.  Rank and exact energy degeneracy are
therefore insufficient to determine the curvature point process.

The curvature form factor provides a second comparison.  Over the registered
finite-$D$ window, the physical curve follows the Jacobi reference ramp within
the uncertainty obtained from complete reference draws.  At longer scales the
difference need not remain small.  Figure~\ref{fig:spectral}(d) and
Fig.~\ref{fig:hierarchy}(d) show the surviving form-factor and number-variance
residuals.  We describe this as finite-rank ramp agreement with long-range
memory.  The statement is tied to the displayed rank, normalization, and
tangent ensemble; it is not an asymptotic or universal claim.

\subsection{The order in which correlations appear}

To separate local and mesoscopic scales, the registered interpolation varies
the strength $g$ of the geometric scrambling while leaving the comparison
protocol fixed.  Figure~\ref{fig:hierarchy}(a) plots the curvature-form-factor
residual over Fourier scale and $g$.  The residual becomes small first in a
restricted interval of Fourier scales.  This flow is not uniform: a local
spacing criterion becomes compatible with the finite-$D$ reference before
the first registered ramp window appears, as summarized in
Fig.~\ref{fig:hierarchy}(b).

The ordering is physically useful.  Nearest-neighbor repulsion tests only the
shortest correlations of the curvature eigenvalues.  A ramp requires
correlations across a wider part of the point process, and number variance
integrates correlations over a still longer window.  Agreement of the first
statistic thus need not force agreement of the other two.  In the physical
panel, the registered ramp residual becomes compatible with the finite-$D$
reference over the shaded window in Fig.~\ref{fig:hierarchy}(c), while the
number-variance difference grows beyond the marked local scale in panel (d).

The uncertainty unit changes with neither Fourier time nor window length.  A
complete physical orbit or a complete ensemble draw is one realization; the
many eigenvalues and time samples within it are correlated observations.  In
particular, treating every curvature eigenvalue as an independent sample
would artificially narrow the bands in Figs.~\ref{fig:spectral} and
\ref{fig:hierarchy}.  The complete physical panel is the statistical unit
throughout this local-to-global comparison.

This hierarchy supplies the precise sense in which the geometry is
Jacobi-like.  It refers to local repulsion and a registered finite-rank ramp,
not to equality of the entire joint distribution with an invariant ensemble.
The long-range residual is evidence that the response retains information
about its physical construction after the local spectrum has become
compatible with the reference.

\subsection{Intrafiber and off-fiber interventions}

The separation between energy and geometry can also be imposed rather than
inferred.  Figure~\ref{fig:channels} compares two interventions on the same
block decomposition.  The intrafiber interpolation changes $PHP$.  Its level
statistics approach the registered GUE interval, but the projector is held
fixed.  The norm $\lVert P_\alpha-P_0\rVert_F$ and the curvature-spectrum
error remain at numerical zero, so this operation changes the spectral
channel without moving the quantum geometry.

The complementary interpolation changes the off-fiber block
$P(\partial H)Q+\mathrm{H.c.}$ while leaving the intrafiber intervention
absent.  The curvature gap statistic rapidly enters its Haar--Jacobi
reference interval.  Thus intrafiber spectral scrambling and off-fiber
projector motion are independent response channels.  Here
\emph{independent} means independently controllable in the block
intervention; it does not mean that estimators measured from one physical
panel are independent random variables.

The four quadrants in Fig.~\ref{fig:channels}(d) are therefore all logically
possible: neither channel may be random-matrix-like, either one can be
compatible while the other remains structured, or both can be compatible.
An exactly degenerate parent begins with a silent intrafiber spectrum, but it
is not forced into the geometrically structured quadrant.  This constructed
control also prevents us from interpreting agreement in the curvature
spectrum as indirect evidence that an unmeasured energy spectrum has become
chaotic.

\subsection{Historical Wick residual and the complete-cumulant boundary}

We next turn from curvature eigenvalues to response matrix elements.  The
Kapit--Mueller calculations use particle numbers $N=3,4,5$, with the audited
quantities collected here to expose the finite-size range:
\begin{table}[b]
\caption{\label{tab:laughlin-matrix}
Kapit--Mueller zero-mode ranks, external gaps, and the historical
four-channel residual.  Every number is supplied by the generated v13
registry macros.  The last two columns refer to the restricted separable
proper-complex two-pairing null, not to the corrected complete cumulant.}
\begin{ruledtabular}
\begin{tabular}{ccccc}
$N$ & $D$ & $\Delta$ & $R_4^{\rm 2pair}$ &
$R_4^{\rm 2pair}-R_{4,\mathrm{Wick}}^{\rm 2pair}$ \\
\hline
$3$ & \LaughlinNThreeRank & \LaughlinNThreeExternalGap
& \LaughlinNThreePhysicalTwoPairRFour & \LaughlinNThreeFourthMomentExcess \\
$4$ & \LaughlinNFourRank & \LaughlinNFourExternalGap
& \LaughlinNFourPhysicalTwoPairRFour & \LaughlinNFourFourthMomentExcess \\
$5$ & \LaughlinNFiveRank & \LaughlinNFiveExternalGap
& \LaughlinNFivePhysicalTwoPairRFour & \LaughlinNFiveFourthMomentExcess \\
\end{tabular}
\end{ruledtabular}
\end{table}
The corresponding Fock-space dimensions are
\LaughlinNThreeBasisDimension, \LaughlinNFourBasisDimension, and
\LaughlinNFiveBasisDimension.  Kernel count, gap, gauge invariance, leakage,
and resolvent gates pass for each row before the matrix statistic is formed.

The statistic in Table~\ref{tab:laughlin-matrix} and
Figs.~\ref{fig:wick}(a,b) is historical.  Its Gaussian null assumes a
separable covariance and a proper complex response, for which the
pseudocovariance pairing is set to zero.  It consequently subtracts two Wick
pairings.  The physical local-density panels and the structured Fourier panel
lie above that restricted null over the registered sizes.  Their residuals
decrease with $N$, as do the generated excesses in the last column.

This is a valid rejection of the stated restricted null, but it is not the
corrected complete cumulant of Eq.~\eqref{eq:wickresidual}.  A positive
historical two-pairing residual can be produced by nonseparable two-point
covariance or nonzero pseudocovariance without any connected four-point
cumulant.  We therefore do not call Figs.~\ref{fig:wick}(a,b) a detection of
non-Gaussian closure failure.  A corrected claim would require the full entry
covariance, all three complex Wick pairings, and uncertainty over declared
independent units.  The Laughlin row in Table~\ref{tab:model-evidence} records
that this complete-covariance gate has not been tested.

The physical error bars in Fig.~\ref{fig:wick}(a) use
\RawMomentResidualRealizationCount\ complete local-density panels.  One entire panel,
not one entry of $X_a$ or one pair of indices, is resampled.  The Fourier
panel is shown as a structured comparison and is not pooled with the local
panels to manufacture an uncertainty estimate.

\subsection{Parent dependence within two-body clustering}

The continuum lowest-Landau-level calculation preserves the same two-body
Laughlin constraint mechanism but changes the microscopic orbitals and the
realization of the physical tangents.  Figure~\ref{fig:wick}(c) displays
guiding-density, magnetic-bond, and magnetic-current panels separately.  For
the size-indexed registry branch, the exported connected excesses are
\ContinuumLaughlinNThreeConnectedExcess,
\ContinuumLaughlinNFourConnectedExcess, and
\ContinuumLaughlinNFiveConnectedExcess.  The figure, rather than those three
representative macros, resolves the full operator-class dependence and its
complete-panel bootstrap interval.

The comparison between the Kapit--Mueller and continuum LLL parents is
deliberately not a pooled test.  Both possess a positive factorized
$V_0$ parent and the same exact constraint-transport identity, but their
complement spectra and local generator maps differ.  Agreement would have
supported a parent-insensitive scalar law within this protection mechanism;
the registered outcome instead selects parent dependence.

Specifically, the sealed common-plateau prediction was evaluated without
refitting on the three continuum tangent classes and three registered sizes.
Figure~\ref{fig:wick}(d) reports that zero of nine registered predictions
cover their corresponding observed points.  The error bars there are
complete-panel observation-bootstrap intervals; the registered prediction is
not adjusted after seeing the continuum outcomes.

The no-refit failure has a narrow meaning.  It rejects that particular scalar
law as a simultaneous description of the displayed operator classes and
sizes.  It does not imply the absence of finite-size decay: all three
continuum classes decrease over the shown range.  It also does not show that
no refitted function could describe the data, that the two microscopic
parents occupy different phases, or that two-body clustering lacks a common
local curvature regime.  The result is parent and operator dependence of the
tested four-channel observable at finite rank.  It supplies no universality
or asymptotic conclusion.
